%%%%%%%%%%%%%%%%%%%%%%%%%%%%%%%%%%%%%%%%%%%%%%%%%%%%%%%%%%%%%%%%%%%%%%
% Lớp 12 | Nguyên hàm, tích phân | Tích phân
% Dạng 2D4?2-3: Tích phân cơ bản hàm lượng giác
% Số câu: 19  (sinh tự động bởi tach_cau.py; từ nay thầy sửa trực tiếp file này)
%%%%%%%%%%%%%%%%%%%%%%%%%%%%%%%%%%%%%%%%%%%%%%%%%%%%%%%%%%%%%%%%%%%%%%

%%%%%%%%%%%%%%%% TRẮC NGHIỆM 4 PHƯƠNG ÁN %%%%%%%%%%%%%%%%

% nguồn: Phan1_Lop_12_Dai_so_Chuong_4
\begin{ex}%[2D4N2-3]%[ma:5dafa82a]
	Tích  phân $\displaystyle \int_0^\pi 2\sin x \mathrm{\,d}x $ bằng
	\choice 
		 {\True $4$}
		{$-4$} 
		 {$2$}
		 {$-2$}
	 \loigiai{
	Ta có $\displaystyle \int_0^\pi  2\sin x \mathrm{\,d}x=-2\cos x \big|_0^\pi=-2(\cos \pi -\cos 0)=4.$}
\end{ex}

% nguồn: Phan1_Lop_12_Dai_so_Chuong_4
\begin{ex}%[2D4N2-3]%[ma:fdd422fc]
	Tích phân $I=\displaystyle\int_0^{\tfrac{\pi}{3}} \cos x\mathrm{\,d}x$  bằng
	\choice 
		 {\True $\dfrac{\sqrt{3}}{2}$}
		 {$-\dfrac{\sqrt{3}}{2}$}
		{$\dfrac{1}{2}$} 
		{$-\dfrac{1}{2}$} 
	\loigiai{ 
		Ta có $I=\displaystyle\int_0^{\tfrac{\pi}{3}}  \cos x\mathrm{\,d}x=\sin x\Bigg|_0^{\tfrac{\pi}{3}}=\dfrac{\sqrt{3}}{2}$.
	 }
\end{ex}

% nguồn: Phan1_Lop_12_Dai_so_Chuong_4
\begin{ex}%[2D4H2-3]%[ma:81afc1b5]
	Biết tích phân $\displaystyle \int_{\tfrac{\pi}{4}}^{\tfrac{2\pi}{3}}\dfrac{3}{\sin^2 x}\mathrm{\,d}x=a+b\sqrt{3}$ ($a$, $b \in  \mathbb{Z}$). Tính $a^2+b^2$.
	 \choice
		{\True  $10$}
		{$4$} 
		 {$\dfrac{9}{2}$}
		{$2$} 
	\loigiai{Ta có $\displaystyle \int_{\tfrac{\pi}{4}}^{\tfrac{2\pi}{3}}\dfrac{3}{\sin^2 x}\mathrm{\,d}x= -3 \cot x  \Biggr|_{\tfrac{\pi}{4}}^{\tfrac{2\pi}{3}}=-3 \left(\cot \dfrac{2 \pi}{3} -\cot \dfrac{\pi}{4}\right)=-3\left(-\dfrac{\sqrt{3}}{3}-1\right)=3+\sqrt{3}$.\\
	Vậy  $a=3$, $b=1$ nên $a^2+b^2=3^2+1^2=10$.}
\end{ex}

% nguồn: Phan1_Lop_12_Dai_so_Chuong_4
\begin{ex}%[2D4H2-3]%[ma:cfdeba3f]
	Giá trị của $\displaystyle\int_0^{\tfrac\pi2}{\sin x\mathrm{\,d}x}$ bằng 
	 \choice
		{$-1$} 
		 {$\dfrac\pi2$}
		{$0$} 
		{\True  $1$}
	\loigiai{ 
		Ta có \[\displaystyle\int_0^{\tfrac\pi2}{\sin x\mathrm{\,d}x}=\left(-\cos x\right) \Bigg|_0^{\tfrac\pi2}=-\cos\dfrac\pi2+\cos 0=1.\] 
	 }
\end{ex}

% nguồn: Phan1_Lop_12_Dai_so_Chuong_4
\begin{ex}%[2D4H2-3]%[ma:5f4b7e19]
	Cho hàm số $f(x)=\heva{&2x^2-x&&\text{ khi }x<0\\&\sin x&&\text{ khi }x\geq 0}$. Tính tích phân $\displaystyle\int_{-1}^{\pi}f(x)\mathrm{\,d}x$. 
	 \choice
		{$\dfrac{13}{6}$} 
		{$\dfrac{5}{6}$} 
		 {$-\dfrac{5}{6}$}
		{\True $\dfrac{19}{6}$} 
	 \loigiai{
		Ta có \[I=\int_{-1}^0f(x)\mathrm{\,d}x+\int_0^{\pi}f(x)\mathrm{\,d}x=\int_{-1}^0\left(2x^2-x\right)\mathrm{\,d}x+\int_0^{\pi}\sin x\mathrm{\,d}x=\dfrac{7}{6}+2=\dfrac{19}{6}.\] 
	} 
\end{ex}

% nguồn: Phan1_Lop_12_Dai_so_Chuong_4
\begin{ex}%[2D4H2-3]%[ma:805ff1ae]
	Cho$\displaystyle\int_0^{\tfrac\pi2}{f(x)}\mathrm{\,d}x=4$. Khi đó $\displaystyle\int_0^{\tfrac\pi2}{\left[2f(x)+\sin x\right]}\mathrm{\,d}x$ bằng 
	\choice 
		{$8+\dfrac\pi2$} 
		{$4+\pi$} 
		{\True  $9$}
		{$7$} 
	 \loigiai{
		Ta có \[\displaystyle\int_0^{\tfrac\pi2}{\left[2f(x)+\sin  x\right]}\mathrm{\,d}x=2\displaystyle\int_0^{\tfrac\pi2}{f(x)}\mathrm{\,d}x+\displaystyle\int_0^{\tfrac\pi2}{\sin x}\mathrm{\,d}x=2\cdot4+1=9.\]
	} 
\end{ex}

% nguồn: Phan1_Lop_12_Dai_so_Chuong_4
\begin{ex}%[2D4V2-3]%[ma:5f12506f]
	Cho $\displaystyle\int_{\tfrac{\pi }6}^{\tfrac{\pi }4}{\cos 4x\cos x\mathrm{\,d}x}=\dfrac{\sqrt2}a+\dfrac bc$ với $ a$, $b$, $c$ là  các số nguyên, $ c<0$ và $\dfrac bc$ tối giản. Tổng $ a+b+c$ bằng
	 \choice
		{$-77$} 
		{$103$} 
		{\True  $-17$}
		 {$43$}
	 \loigiai{
		 Ta có
		 \allowdisplaybreaks
		\begin{eqnarray*}
			\displaystyle\int_{\tfrac{\pi }6}^{\tfrac{\pi }4}{\cos  4x\cos x\mathrm{\,d}x}&=&\dfrac12\,\displaystyle\int_{\tfrac{\pi }6}^{\tfrac{\pi }4}{\left(\cos5x+\cos3x\right)\mathrm{\,d}x}\\
			&=&\dfrac12\left(\dfrac15\sin 5x+\dfrac13\sin3x\right) \Bigg|_{\tfrac{\pi  }6}^{\tfrac{\pi }4}\\
			&=&\dfrac{\sqrt2}{30}-\dfrac{13}{60}. 
		\end{eqnarray*}
		 Suy ra
		\[\heva{ 
		 &a=30\\
		 &b=13\\
		&c=-60\\ 
		}\Rightarrow  a+b+c=30+13-60=-17.\]
	 }
\end{ex}

% nguồn: Phan1_Lop_12_Dai_so_Chuong_4
\begin{ex}%[2D4V2-3]%[ma:76a386d5]
	Cho biết $\displaystyle\int_0^{\tfrac\pi2}{\left(4-\sin x\right)\mathrm{\,d}x}=a\pi+b$, với $ a,b$ là các số nguyên.  Giá trị của biểu thức $ a+b$ bằng
	\choice 
		{\True  $1$}
		{$-4$} 
		 {$6$}
		 {$3$}
	\loigiai{ 
		Ta có 
		\[\displaystyle\int_0^{\tfrac\pi2}{\left(4-\sin x\right)\mathrm{\,d}x}=\left(4x+\cos  x\right) \Bigg|_0^{\tfrac\pi2}=2\pi-1\Rightarrow a=2,b=-1.\]
	 Vậy $ a+b=1.$}
\end{ex}

% nguồn: Phan1_Lop_12_Dai_so_Chuong_4
\begin{ex}%[2D4V2-3]%[ma:2f0eeec2]
	Biết $\displaystyle\int_0^{\tfrac{\pi }4}{\dfrac1{{{\sin }^2}x.{{\cos }^2}x}\mathrm{\,d}x}=\dfrac{a\sqrt3}b$ $\left(a,b\in\mathbb{Z}\right)$. Tính $P=\dfrac{a-2b}b$ 
	\choice 
		 {$P=\dfrac43$}
		{\True $P=-\dfrac43$} 
		{$P=-\dfrac23$} 
		{$P=\dfrac23$} 
	\loigiai{Ta có 
		\begin{align*}
			\displaystyle\int_{\tfrac{\pi }4}^{\tfrac{\pi }3}{\dfrac1{\sin^2x\cdot{{\cos }^2}x}\mathrm{\,d}x}&=\displaystyle\int_{\tfrac{\pi }4}^{\tfrac{\pi }3}{\dfrac{{{\sin }^2}x+{{\cos }^2}x}{{{\sin }^2}x\cdot{{\cos  }^2}x}\mathrm{\,d}x}\\
			&=\displaystyle\int_{\tfrac{\pi }4}^{\tfrac{\pi }3}{\left(\dfrac1{{{\cos }^2}x} \text{  +}\dfrac1{{{\sin }^2}x}\right)\mathrm{\,d}x}\\
			&=\left(\tan x-\cot  x\right) \Bigg|_{\tfrac{\pi }4}^{\tfrac{\pi }3}\\
			&=\dfrac{2\sqrt3}3. 
		\end{align*}
		Với  $a=2$, $b=3$, vậy $P=\dfrac{2-2\cdot3}3=-\dfrac43.$
	} 
\end{ex}

% nguồn: Phan1_Lop_12_Dai_so_Chuong_4
\begin{ex}%[2D4V2-3]%[ma:55f00d6a]
	Cho hàm số $f(x)=\heva{&3x^2+a&&\text{ khi }x\geq 0\\&\sin 2x+\cos x-b&&\text{ khi }x<0}$ liên tục và có nguyên hàm $F$ liên tục trên $\mathbb{R}$ thỏa mãn $F(2)-F\left(-\dfrac{\pi}{2}\right)=10$ với $a,b$ là tham số thực. Giá trị $\displaystyle\int_{-\tfrac{\pi}{6}}^1f(x)\mathrm{\,d}x=\dfrac{m}{4}$, với $m$ là  một số thực. Khi đó $2m$ bằng
	 \choice
		 {$\dfrac{9}{2}$}
		{$\dfrac{9}{4}$} 
		 {$9$}
		{\True  $18$}
	\loigiai{ 
		Hàm số $f(x)$ liên tục  trên $\mathbb{R}$ nên $a=1-b$.\\
		 Mặt khác
		\allowdisplaybreaks\begin{eqnarray*}
		 &&F(2)-F\left(-\dfrac{\pi}{2}\right)=10\\
		&\Leftrightarrow&  \int_{-\tfrac{\pi}{2}}^2f(x)\mathrm{\,d}x=10\\
		 &\Leftrightarrow & \int_{-\tfrac{\pi}{2}}^0(\sin 2x+\cos x-b)\mathrm{\,d}x+\int_0^2\left(3x^2+a\right)\mathrm{\,d}x=10\\
		&\Leftrightarrow  &\left(-\dfrac{1}{2}\cos 2x+\sin x-bx\right)\bigg|_{-\tfrac{\pi}{2}}^0+\left(x^3+ax\right)\bigg|_0^2=10\\
		&\Leftrightarrow& 2a-\dfrac{b\pi}{2}=2 
	\end{eqnarray*}
	Do đó, ta  có hệ
	\[\heva{&a+b=1\\ 
	&2a-\dfrac{b\pi}{2}=2}\Leftrightarrow  \heva{&a=1\\
	&b=0.}\] 
	 Suy ra \[f(x)=\heva{&3x^2+1&&\text{ khi }x\geq 0\\
	&\sin 2x+\cos  x&&\text{ khi }x<0.}\]
	 Ta được
	 \allowdisplaybreaks
	\begin{eqnarray*}
		\int_{-\tfrac{\pi}{6}}^1f(x)\mathrm{\,d}x&=&\int_{-\tfrac{\pi}{6}}^0(\sin 2x+\cos  x)\mathrm{\,d}x+\int_0^1\left(3x^2+1\right)\mathrm{\,d}x\\
		&=&\left(-\dfrac{1}{2}\cos 2x+\sin x\right)\bigg|_{-\tfrac{\pi}{6}}^0+\left(x^3+x\right)\bigg|_0^1\\ 
		 &=&\dfrac{1}{4}+2=\dfrac{9}{4}.
	\end{eqnarray*}
 Vậy $m=9$ và $2m=18$.}
\end{ex}

% nguồn: Phan1_Lop_12_Dai_so_Chuong_4
\begin{ex}%[2D4V2-3]%[ma:ceab4f5f]
	Tính tích  phân $I=\displaystyle\int_0^{\tfrac{\pi }4}{\sin3x\sin x\mathrm{\,d}x}$.
	 \choice
		 {$\dfrac12$}
		{$0$} 
		 {$-\dfrac12$}
		 {\True $\dfrac14$}
	\loigiai{ 
		Ta  có
		 \allowdisplaybreaks
		\begin{eqnarray*}
			I=\displaystyle\int_0^{\tfrac{\pi }4}{\sin3x\sin  x\mathrm{\,d}x}&=&\dfrac12\displaystyle\int_0^{\tfrac{\pi }4}{\left(\cos2x-\cos 4x\right)\mathrm{\,d}x}\\
			&=&\dfrac12\left(\dfrac12\sin2x-\dfrac14\sin4x\right)  \Bigg|_0^{\tfrac{\pi }4}\\
			 &=&\dfrac12\left(\dfrac12-0\right)\\
			&=&\dfrac14. 
		\end{eqnarray*}
	} 
\end{ex}

% nguồn: Phan1_Lop_12_Dai_so_Chuong_4
\begin{ex}%[2D4V2-3]%[ma:704645f0]
	Cho hàm  số $f(x)=\heva{&4x-\sqrt{4x+9}&&\text{ khi }x>0\\&4a+\tan^2x&&\text{ khi }x\leq 0}$ và $I=\int_{-\tfrac{\pi}{4}}^4f(x)\mathrm{\,d}x=\dfrac{50}{3}$. Tính $a$.
	 \choice
		{$a=1$} 
		 {$a=\dfrac{1}{2}$}
		{$a=\dfrac{3}{4}$} 
		{\True  $a=\dfrac{1}{4}$}
	 \loigiai{
		Ta có \allowdisplaybreaks\begin{eqnarray*}
		 I&=&\int_{-\tfrac{\pi}{4}}^4f(x)\mathrm{\,d}x\\
		&=&\int_{-\tfrac{\pi}{4}}^0f(x)\mathrm{\,d}x+\int_0^4f(x)\mathrm{\,d}x\\ 
		 &=&\int_{-\tfrac{\pi}{4}}^0\left(4a+\tan^2x\right)\mathrm{\,d}x+\int_0^4\left(4x-\sqrt{4x+9}\right)\mathrm{\,d}x\\
		&=&\int_{-\tfrac{\pi}{4}}^0(4a-1)\mathrm{\,d}x+\int_{-\tfrac{\pi}{4}}^0\left(1+\tan^2x\right)\mathrm{\,d}x+\int_0^4\left(4x-\sqrt{4x+9}\right)\mathrm{\,d}x\\ 
		&=&(4a-1)\bigg|_{-\tfrac{\pi}{4}}^0+\tan  x\bigg|_{-\tfrac{\pi}{4}}^0+\left(2x^2-\dfrac{\sqrt{(4x+9)^3}}{6}\right)\bigg|_0^4\\
		&=&(4a-1)\pi+1+\dfrac{47}{3}. 
	\end{eqnarray*}
	Do đó  \[I=\dfrac{50}{3}\Leftrightarrow(4a-1)\pi+1+\dfrac{47}{3}=\dfrac{50}{3}\Leftrightarrow a=\dfrac{1}{4}.\]
} 
\end{ex}

% nguồn: Phan1_Lop_12_Dai_so_Chuong_4
\begin{ex}%[2D4V2-3]%[ma:13f280f9]
	Biết $\displaystyle\int_{-\tfrac{\pi}{2}}^{\tfrac{\pi}{6}}|\sin x|\mathrm{\,d}x=a-\sqrt{b}$ ($a$, $b \in \mathbb{Q})$. Khi đó $a+4b$ bằng 
	 \choice
		{\True $5$} 
		 {$8$}
		 {$10$}
		{$7$} 
	 \loigiai{
		Ta có \begin{eqnarray*}
		 \int_{-\tfrac{\pi}{2}}^{\tfrac{\pi}{6}}|\sin x|\mathrm{\,d}x&=&-\int_{-\tfrac{\pi}{2}}^0\sin x \mathrm{\,d}x+\int_0^{\tfrac{\pi}{6}} \sin x \mathrm{\,d}x\\
		&=&(\cos x)\bigg|_{-\tfrac{\pi}{2}}^0-(\cos x)\bigg|_0  ^{\tfrac{\pi}{6}}\\
		 &=&2-\dfrac{\sqrt{3}}{2}\\
		&=&2-\sqrt{\dfrac{3}{4}}. 
	\end{eqnarray*}
	Suy ra $a=2$,  $b=\dfrac{3}{4}$ nên $a+4b=5$.
 }
\end{ex}

% nguồn: Phan1_Lop_12_Dai_so_Chuong_4
\begin{ex}%[2D4V2-3]%[ma:6b11f399]
	Cho hàm số $f(x)=\heva{&2x+a&&\text{ khi }x\geq 1\\ &3x^2+b&&\text{ khi }x<1}$  thỏa mãn $\displaystyle\int_0^2f(x)\mathrm{\,d}x=13$. Tính giá trị biểu thức $T=a+b-ab$?
	\choice 
		 {\True $T=-11$}
		{$T=-5$} 
		{$T=1$} 
		{$T=-1$} 
	\loigiai{ 
		Để tồn tại $\displaystyle\int_0^2f(x)\mathrm{\,d}x$ thì $f(x)$  liên tục trên đoạn $[0;2]$, suy ra $f(x)$ liên tục tại $x=1$.\\
		Khi  đó
		\[\lim\limits_{x\to 1^{+}}f(x)=\lim\limits_{x\to 1^{-}}f(x)=f(1)\Leftrightarrow  a+2=b+3\Leftrightarrow a=b+1.\quad(1)\]
		 Ta có
		\allowdisplaybreaks 
		\begin{eqnarray*}
			\int_0^2f(x)\mathrm{\,d}x&=&\int_0^1f(x)\mathrm{\,d}x+\int_1^2f(x)\mathrm{\,d}x\\ 
			 &=&\int_0^1\left(3x^2+b\right)\mathrm{\,d}x+\int_1^2(2x+a)\mathrm{\,d}x\\
			&=&\left(x^3+bx\right)\bigg|_0^1+\left(x^2+ax\right)\bigg|_1^2\\ 
			 &=&a+b+4.
		\end{eqnarray*}
		 Mà $\displaystyle\int_0^2f(x)\mathrm{\,d}x=13\Rightarrow a+b=9$. $(2)$\\
		 Từ $(1)$ và $(2)$ suy ra $a=5$, $b=4$.\\
		 Vậy $T=a+b-ab=-11$.
	 }
\end{ex}

% nguồn: Phan1_Lop_12_Dai_so_Chuong_4
\begin{ex}%[2D4C2-3]%[ma:79da5039]
	Cho hàm số $f(x)=\heva{&\sin x+2&&\text{ khi }x\geq 0\\&2\cos^2x&&\text{ khi }x<0}$. Giả sử $F$ là nguyên hàm của $f$ trên $\mathbb{R}$ thỏa  mãn $F\left(\dfrac{\pi}{3}\right)=\pi$. Giá trị của $F\left(-\dfrac{\pi}{6}\right)+2F\left(\dfrac{\pi}{6}\right)-F\left(\dfrac{\pi}{4}\right)$ bằng
	\choice 
		{$\dfrac{1+\sqrt{3}+3\pi}{2}$} 
		 {\True $\dfrac{1}{2}-\dfrac{\sqrt{3}}{2}+\dfrac{13\pi}{12}$}
		{$\dfrac{13\pi+6-6\sqrt{3}}{12}$} 
		 {$\dfrac{1}{2}-\dfrac{\sqrt{3}}{2}+\dfrac{\pi}{12}$}
	 \loigiai{
		Ta có \allowdisplaybreaks\begin{eqnarray*}
		 &&F\left(-\dfrac{\pi}{6}\right)+2F\left(\dfrac{\pi}{6}\right)-F\left(\dfrac{\pi}{4}\right)\\
		 &=&F\left(-\dfrac{\pi}{6}\right)-F(0)+F(0)-F\left(\dfrac{\pi}{6}\right)+3F\left(\dfrac{\pi}{6}\right)-3F\left(\dfrac{\pi}{4}\right)+2F\left(\dfrac{\pi}{4}\right)-2F\left(\dfrac{\pi}{3}\right)+2F\left(\dfrac{\pi}{3}\right)\\
		&=&-\int_{-\tfrac{\pi}{6}}^0f(x)\mathrm{\,d}x-\int_0^{\tfrac{\pi}{6}}f(x)\mathrm{\,d}x-3\int_{\tfrac{\pi}{6}}^{\tfrac{\pi}{4}}f(x)\mathrm{\,d}x-2\int_{\tfrac{\pi}{4}}^{\tfrac{\pi}{3}}f(x)\mathrm{\,d}x+2F\left(\dfrac{\pi}{3}\right)\\ 
		 &=&-\int_{-\tfrac{\pi}{6}}^0(\sin x+2)\mathrm{\,d}x-\int_0^{\tfrac{\pi}{6}}2\cos^2x\mathrm{\,d}x-3\int_{\tfrac{\pi}{6}}^{\tfrac{\pi}{4}}2\cos^2x\mathrm{\,d}x-2\int_{\tfrac{\pi}{4}}^{\tfrac{\pi}{3}}2\cos^2x\mathrm{\,d}x+2F\left(\dfrac{\pi}{3}\right)\\
		&=&\int_{-\tfrac{\pi}{6}}^0(-\sin x-2)\mathrm{\,d}x-\int_0^{\tfrac{\pi}{6}}(\cos 2x+1)\mathrm{\,d}x-3\int_{\tfrac{\pi}{6}}^{\tfrac{\pi}{4}}(\cos  2x+1)\mathrm{\,d}x-2\int_{\tfrac{\pi}{4}}^{\tfrac{\pi}{3}}(\cos 2x+1)\mathrm{\,d}x+2F\left(\dfrac{\pi}{3}\right)\\
		&=&(\cos x-2x)\bigg|_{-\tfrac{\pi}{6}}^0-\left(\dfrac{1}{2}\sin 2x+x\right)\bigg|_0^{\tfrac{\pi}{6}}-3\left(\dfrac{1}{2}\sin  2x+x\right)\bigg|_{\tfrac{\pi}{6}}^{\tfrac{\pi}{4}}-2\left(\dfrac{1}{2}\sin 2x+x\right)\bigg|_{\tfrac{\pi}{4}}^{\tfrac{\pi}{3}}+2F\left(\dfrac{\pi}{3}\right)\\
		 &=&1-\dfrac{\sqrt{3}}{2}-\dfrac{\pi}{3}-\left(\dfrac{\sqrt{3}}{4}+\dfrac{\pi}{6}\right)-3\left(\dfrac{1}{2}-\dfrac{\sqrt{3}}{4}+\dfrac{\pi}{12}\right)-2\left(\dfrac{\sqrt{3}}{4}-\dfrac{1}{2}+\dfrac{\pi}{12}\right)+2\pi\\
		&=&\dfrac{1}{2}-\dfrac{\sqrt{3}}{2}+\dfrac{13\pi}{12}. 
	\end{eqnarray*}
} 
\end{ex}


%%%%%%%%%%%%%%%% ĐÚNG/SAI %%%%%%%%%%%%%%%%

% nguồn: Phan1_Lop_12_Dai_so_Chuong_4
\begin{ex}%[2D4H2-3]%[ma:136aaed2]
	 Cho hàm số $f(x)=x-\sin x$.
	 \choiceTF
		{\True Gọi $F(x)$ là một nguyên hàm của $f(x)$. Khi đó $F'(x)=f(x)$} 
		{Hàm số $F(x)=\dfrac{x^2}{2}-\cos x+C$ là họ các nguyên hàm của hàm số  $f(x)$}
		{\True $\displaystyle\int_0^{\tfrac{\pi}{2}}f(x)\mathrm{\,d}x = \int_0^{\tfrac{\pi}{2}}x\mathrm{\,d}x -  \int_0^{\tfrac{\pi}{2}}\sin x\mathrm{\,d}x$}
		{\True $\displaystyle\int_0^{\tfrac{\pi}{2}}f(x)\mathrm{\,d}x = \dfrac{\pi^2}{8}  - 1$}
	\loigiai{ 
		\begin{itemchoice}
			\itemch \textbf{Đúng}.  Theo định nghĩa, nếu $F(x)$ là một nguyên hàm của $f(x)$. Khi đó $F'(x)=f(x)$.
			\itemch \textbf{Sai}. Vì $\displaystyle\int f(x)\mathrm{\,d}x = \int (x - \sin x)\mathrm{\,d}x = \dfrac{x^2}{2} + \cos x +  C$.\\
			Do đó, họ các nguyên hàm của hàm số $f(x)$ là $\dfrac{x^2}{2} + \cos x  + C$.
			\itemch \textbf{Đúng}. Ta có $\displaystyle\int_0^{\tfrac{\pi}{2}} f(x)\mathrm{\,d}x  = \int_0^{\tfrac{\pi}{2}} (x - \sin x)\mathrm{\,d}x = \int_0^{\tfrac{\pi}{2}}x\mathrm{\,d}x - \int_0^{\tfrac{\pi}{2}}\sin x\mathrm{\,d}x$.
			\itemch \textbf{Đúng}.  Ta có
			\allowdisplaybreaks 
			\begin{eqnarray*}
				\int_0^{\tfrac{\pi}{2}} f(x)\mathrm{\,d}x = \int_0^{\tfrac{\pi}{2}} (x - \sin x)\mathrm{\,d}x &=& \int_0^{\tfrac{\pi}{2}}x\mathrm{\,d}x - \int_0^{\tfrac{\pi}{2}}\sin x\mathrm{\,d}x\\ 
				&=& \dfrac{x^2}{2}\Big|_{0}^{\tfrac{\pi}{2}} + \cos x\Big|_{0}^{\tfrac{\pi}{2}}\\ 
				&=& \dfrac{\left(\dfrac{\pi}{2}\right)^2}{2} - \dfrac{0^2}{2} + \cos\dfrac{\pi}{2} -  \cos 0\\
				&=&  \dfrac{\pi^2}{8} - 1.
			\end{eqnarray*}
		\end{itemchoice}
	 }
\end{ex}

% nguồn: Phan1_Lop_12_Dai_so_Chuong_4
\begin{ex}%[2D4H2-3]%[ma:5211af9f]
	Cho $f(x)=\sin  x$. Khi đó
	 \choiceTF
		{\True $\displaystyle\int_0^{\tfrac{\pi}{4}}  f(x)\mathrm{\,d}x=1-\dfrac{\sqrt{2}}{2}$}
		{\True $\displaystyle\int_0^{\pi} f\left( 2x-\dfrac{\pi}{3} \right)\mathrm{\,d}x=0$} 
		{$\displaystyle\int_0^{\tfrac{\pi}{2}}\big[2x-f(x)\big]\mathrm{\,d}x=\dfrac{\pi^2}{4}+1$} 
		{\True $\displaystyle\int_0^{\tfrac{\pi}{2}}  xf(2x)\mathrm{\,d}x=\dfrac{\pi}{4}$}
	\loigiai{ 
		\begin{itemchoice}
			\itemch Ta có $\displaystyle\int_0^{\tfrac{\pi}{4}}  f(x)\mathrm{\,d}x=\displaystyle\int_0^{\tfrac{\pi}{4}} \sin x\mathrm{\,d}x=1-\dfrac{\sqrt{2}}{2}$.
			\itemch Ta có $\displaystyle\int_0^{\pi} f\left( 2x-\dfrac{\pi}{3} \right)\mathrm{\,d}x=\displaystyle\int_0^{\pi}  \sin \left( 2x-\dfrac{\pi}{3} \right)\mathrm{\,d}x=0$.
			\itemch  Ta có $\displaystyle\int_0^{\tfrac{\pi}{2}} \big[2x-f(x)\big]\mathrm{\,d}x=\displaystyle\int_0^{\tfrac{\pi}{2}} 2x\mathrm{\,d}x-\displaystyle\int_0^{\tfrac{\pi}{2}} \sin x\mathrm{\,d}x=\dfrac{\pi^2}{4}-1$.
			\itemch Ta có $\displaystyle\int_0^{\tfrac{\pi}{2}} xf(2x)\mathrm{\,d}x=\displaystyle\int_0^{\tfrac{\pi}{2}} x\cdot\sin  2x\mathrm{\,d}x=\dfrac{\pi}{4}$.
		\end{itemchoice}
	} 
\end{ex}


%%%%%%%%%%%%%%%% TRẢ LỜI NGẮN %%%%%%%%%%%%%%%%

% nguồn: Phan2_Lop_12_Dai_so_Chuong_4
\begin{bt}%[2D4H2-3]%[ma:d5f101df]
	Kết quả của tích  phân $\displaystyle\int_0^{\tfrac{\pi}{4}} \sin x\mathrm{\,d}x = \dfrac{a-\sqrt{b}}{2}$. Tính $a+b$.
	 \shortans{4}
	 \loigiai{
		Ta có $\displaystyle\int_0^{\tfrac{\pi}{4}}  \sin x\mathrm{\,d}x = -\cos x\Bigg|_0^{\tfrac{\pi}{4}} = \left(-\cos\dfrac{\pi}{4}\right)-(-\cos 0)=-\dfrac{\sqrt{2}}{2}-(-1)=\dfrac{2-\sqrt{2}}{2}$.\\
		 Vậy $a+b=2+2=4$.
	} 
\end{bt}


%%%%%%%%%%%%%%%% TỰ LUẬN %%%%%%%%%%%%%%%%

% nguồn: Phan2_Lop_12_Dai_so_Chuong_4
\begin{bt}%[2D4H2-3]%[ma:e8821c43]
	Tính các  tích phân sau:
	\begin{listEX}[2]
		\item[a)] $\displaystyle \int_{\frac{\pi}{6}}^{\frac{\pi}{4}}\dfrac{1}{\sin^2x}\mathrm{\,d}x$; 
		\item  [b)] $\displaystyle \int_{0}^{\frac{\pi}{4}}(1+\tan x)\cos x\mathrm{\,d}x$.
	\end{listEX}
	 \loigiai{
		\begin{enumerate}
			\item $\displaystyle \int_{\frac{\pi}{6}}^{\frac{\pi}{4}}\dfrac{1}{\sin^2x}\mathrm{\,d}x=-\cot  x\bigg|_{\frac{\pi}{6}}^{\frac{\pi}{4}}=-(1-\sqrt{3})=\sqrt{3}-1$;
			\item $\displaystyle \int_{0}^{\frac{\pi}{4}}(1+\tan x)\cos x\mathrm{\,d}x=\displaystyle \int_{0}^{\frac{\pi}{4}}(\cos  x+\sin x)\mathrm{\,d}x=\left(\sin x-\cos x\right)\bigg|_{0}^{\frac{\pi}{4}}=1$.
		\end{enumerate}
	} 
\end{bt}
